% Method Section - Semantic-Structural Decomposition

\subsection{The Limitation of GP Variance}

\begin{proposition}[Relation-Agnostic]
\label{prop:relation_agnostic}
GP-KGE uncertainty (Eq.~\ref{eq:gp_uncertainty}) is relation-agnostic: for any entity $e$ and relations $r_1, r_2$,
\begin{equation}
    U_{\text{GP}}(e, r_1, \cdot) = U_{\text{GP}}(e, r_2, \cdot)
\end{equation}
\end{proposition}

\begin{proof}
Immediate from Eq.~\ref{eq:gp_uncertainty}: the relation $r$ does not appear in the uncertainty computation.
\end{proof}

This limitation is fundamental, not an implementation detail (see Figure~\ref{fig:gp_limitation}).
Making GP variance relation-aware would require learning $|\mathcal{E}| \times |\mathcal{R}| \times d$ parameters---infeasible for typical KG sizes.

\begin{figure}[t]
\centering
\includegraphics[width=0.85\linewidth]{figures/fig5_gp_limitation.pdf}
\vspace{-1em}
\caption{GP variance (right) is constant across relations; coverage (left) is relation-specific.}
\label{fig:gp_limitation}
\end{figure}

\subsection{Structural Uncertainty via Coverage}

To address this limitation, \emph{coverage} serves as a relation-specific uncertainty signal.

\begin{definition}[Coverage]
For entity $e$ and relation $r$:
\begin{equation}
    c(e, r) = \mathbb{1}\left[\exists (h', r, t') \in \mathcal{T} : h' = e \lor t' = e\right]
\end{equation}
\end{definition}

Coverage equals 1 if entity $e$ appears with relation $r$ in training, else 0.

\begin{definition}[Structural Uncertainty]
\begin{equation}
    U_{\text{cov}}(h, r, t) = 2 - c(h, r) - c(t, r)
\end{equation}
\end{definition}

This ranges from 0 (both entities seen with $r$) to 2 (neither seen).

\begin{proposition}[Coverage AUROC under Random Corruption]
\label{prop:coverage_auroc}
Under random tail corruption with relation sparsity $s_r = 1 - |\{e : c(e,r)=1\}|/|\mathcal{E}|$, let $a = P(\text{both covered} | \text{ID})$ and $b = P(\text{exactly one covered} | \text{ID})$. Then:
\begin{equation}
    \text{AUROC}_{\text{cov}} = \frac{1}{2}\left(a(1+s_r) + s_r \cdot b\right)
\end{equation}
\end{proposition}

\begin{proof}[Derivation]
ID uncertainty distribution: $P(U=0)=a$, $P(U=1)=b$, $P(U=2)=1-a-b$.
OOD has head from ID, random tail: $P(U=0)=1-s_r$, $P(U=1)=s_r$.
AUROC follows from $P(U_{\text{ID}} < U_{\text{OOD}}) + \frac{1}{2}P(U_{\text{ID}} = U_{\text{OOD}})$.
Full derivation in Appendix~\ref{app:proof}.
\end{proof}

\subsection{Complementarity of Signals}

The two uncertainty types capture orthogonal information, making neither sufficient alone.

\begin{proposition}[Complementarity]
\label{prop:complementarity}
Coverage and GP variance are not redundant: each correctly identifies OOD samples that the other misses.
\end{proposition}

\begin{proof}
We construct explicit failure cases for each signal:

\textbf{Case 1 (GP fails, Coverage succeeds):} Let entity $e$ appear in 1000 training triples across relations $\{r_1, \ldots, r_{10}\}$, but never with relation $r_{11}$.
GP assigns low variance $\sigma^2_e$ (well-observed entity).
For query $(?, r_{11}, e)$: GP incorrectly signals low uncertainty, while coverage correctly signals high uncertainty since $c(e, r_{11}) = 0$.

\textbf{Case 2 (Coverage fails, GP succeeds):} Let entity $e'$ appear in exactly one training triple $(h, r, e')$.
Coverage assigns $c(e', r) = 1$.
For OOD query $(h', r, e')$ where $e'$ was randomly sampled: Coverage incorrectly signals low uncertainty, while GP correctly signals high uncertainty due to large $\sigma^2_{e'}$.
\end{proof}

\noindent\textbf{Intuition.} GP variance captures \emph{how well an entity is learned} (global property), while coverage captures \emph{whether an entity-relation pair was observed} (local property). A frequently-observed entity can still be queried in novel relational contexts; a rarely-observed entity can still appear in familiar relational contexts.

\subsection{CAGP: Combining Both Signals}

Proposition~\ref{prop:complementarity} motivates combining both signals:

\begin{definition}[CAGP Uncertainty]
\begin{equation}
    U_{\text{CAGP}}(h, r, t) = \alpha \cdot \tilde{U}_{\text{GP}}(h, r, t) + (1-\alpha) \cdot U_{\text{cov}}(h, r, t)
\end{equation}
where $\tilde{U}_{\text{GP}}$ is GP uncertainty normalized to match coverage scale, and $\alpha \in [0,1]$ is a learnable mixing coefficient.
\end{definition}

Implementation details:
\begin{itemize}
    \item Coverage matrix is precomputed from training data (no additional learning)
    \item $\alpha$ is parameterized as $\sigma(\alpha_{\text{logit}})$ to ensure $\alpha \in (0,1)$
    \item Normalization: $\tilde{U}_{\text{GP}} = U_{\text{GP}} \cdot \frac{\mathbb{E}[U_{\text{cov}}]}{\mathbb{E}[U_{\text{GP}}]}$
\end{itemize}

The algorithm adds minimal overhead: one $|\mathcal{E}| \times |\mathcal{R}|$ binary matrix (sparse in practice) and one scalar parameter.
